\documentclass[11pt]{article}
\usepackage[utf8]{inputenc}
\usepackage[T1]{fontenc}
\usepackage[margin=1in]{geometry}
\usepackage{hyperref}
\usepackage{xcolor}
\usepackage{xr}
\usepackage{natbib}

\externaldocument{manuscript}

\DeclareUnicodeCharacter{03B2}{$\beta$}
\DeclareUnicodeCharacter{2013}{--}
\DeclareUnicodeCharacter{2014}{---}
\DeclareUnicodeCharacter{2019}{'}
\DeclareUnicodeCharacter{00B2}{$^2$}

\title{Response to Reviewers}
\author{Zoey Werbin, Colin Averill, Jennifer Bhatnagar, Michael Dietze}
\date{}

\begin{document}
\maketitle

\href{https://docs.google.com/document/d/1SuD6-FiOA-zIBZACd5x7ZjkDE_sT7UG2fEuCYWlSb9Y/edit?usp=sharing}{Link to manuscript} (with supplement)

\bigskip\hrule\bigskip

Referee expertise:

Referee \#1: microbial ecology, computational biology, statistics

Referee \#2: microbial ecology, biogeography

Referee \#3: trait ecology, soil communities, biotic interactions

We thank all the reviewers for their time and detailed, actionable feedback. We hope that this revision satisfactorily addresses the major concerns.

Reviewers' comments:

\section*{Reviewer \#1 (Remarks to the Author)}

\subsection*{Overview}

\textit{This manuscript examines the predictability of soil microbial taxa through time. It considers both the scale over which they can be predicted and the factors that contribute to that predictability, or lack thereof. It finds that in some cases predictability is quite good for up to one year, and is best for microbes grouped at into broad taxonomic groups or ecological groups.}

\textit{This manuscript addresses an exciting and important topic, which I highly commend the authors for tackling. The topic that it addresses is novel, and the paper is very well written. I would love to see a paper on this topic get published. However, there are two major issues that the manuscript needs to address before it will be suitable for publication. These issues will require a major revision:}

\subsection*{Detailed comments}

\textit{(1) Different time series can have different levels of predictability for relatively trivial reasons. For instance, a constant system with no fluctuations or trend can easily be predicted, and a seasonally (but regularly) varying system can likewise be predicted quite simply. More complex systems --- for instance, those exhibiting low amounts of temporal autocorrelation --- are less predictable. Hence, although the manuscript does to some extent investigate seasonal fluctuations, it does not address what differences in predictability really mean. Do they result from some taxa exhibiting linear trends, while others have chaotic trends? Or do they all show evidence of chaotic trends, but with some of those trends having longer horizons of predictability? In essence, the manuscript needs to unpack what differences in predictability mean. This would greatly help in assessing whether the results stem from relatively trivial mechanisms (e.g., some taxa show seasonal trends, while others don't) or more subtle and complex mechanisms, for instance, fundamental differences in the dimensionality of underlying chaotic attractors, perhaps resulting from complicated interspecific interactions. Additional analyses and figures will likely be needed to fully address these questions.}

We agree wholeheartedly with this comment: predictability can arise from a variety of temporal patterns, which can have very different consequences for causal inference and practical management scenarios. We attempted to explore these in a variety of analyses in the supplementary materials: we look at whether predictability is a function of spatial autocorrelation, simple temporal autocorrelation, seasonal trends that can or cannot be explained by environmental predictors, environmental variability (using latitude as a proxy), and sensitivity to environmental predictors with varying spatial and temporal patterns. Despite exploring many potential drivers of predictability in our supplementary analyses, we did not identify a single dominant mechanism. However, our results are most consistent with the interpretation that more predictable taxa and guilds are those whose abundances track measurable environmental gradients more closely. We have added text to the manuscript (lines \ref{rev-1}) clarifying this interpretation and acknowledging that distinguishing among dynamical mechanisms would require experimental or long-term manipulative approaches beyond the scope of this observational dataset.

\bigskip

\textit{(2) As I understand it, the manuscript assesses the predictability of (untransformed) relative abundances of taxa. Because high-throughput sequencing data are compositional, if one taxon undergoes a bloom or population decrease, that can change the relative abundances of the other taxa, even if the actual, non-relative abundances of the latter taxa remain constant or vary differently. Hence, one predictably blooming or crashing taxon could drive similarly predictable trends in the relative abundance of other taxa, even if their populations are actually quite unpredictable: the compositional nature of high-throughput sequence data can alone make the predictabilty of taxa non-independent. The usual way to address this issue would be to predict the ratios of relative abundances of taxa; a well-established approach for both sequence data and copmpositional data generally. However, care must be taken when selecting which ratios to model. The manuscript needs to ensure that the predictability of groups of taxa, or the lack thereof, is not a result of the compositional nature of sequence data and take into account the non-independence that it introduces.}

This is an excellent point, and we agree with the reviewer's thoughtful characterization of composition microbial community data and the caveats associated with its predictability. Our revision now addresses this more explicitly in the main manuscript text. In our previous submission, we conducted thorough comparisons of models built on untransformed relative abundances with models built using ratio transformations, and models built using covariance-aware relative abundances. These results were included at the end of the supplementary material, showing that the major trends in predictability explored by our analyses are robust to methodological decisions about transformations and data modeling. For our broad-scale analysis of soil microbiome trends, we used relative abundances of specific groups for the sake of interpretability, efficiency of model fitting, and generalizability to new datasets; but we acknowledge that there are valid reasons to prefer specific transformations and modeling approaches for more specific use cases. We hope that this rationale is clearer in the updated main manuscript text (lines \ref{rev-2-start}--\ref{rev-2-end}). Specifically, we have added text at lines \ref{rev-3} explaining that while absolute abundances cannot be recovered from amplicon data, the major patterns of predictability we report --- including the taxonomic and functional scaling of forecast skill --- are consistent across untransformed relative abundances, log-ratio transformed data, and covariance-aware models. A detailed comparison of all four observation model formulations --- including the Dirichlet (which explicitly accounts for the sum-to-one constraint), centered log-ratio, truncated normal, and beta regression approaches --- is provided in Appendix S3 (Sections 1--3 and Table 1), along with a figure showing that parameter estimates are qualitatively similar across approaches. We therefore are confident that compositional non-independence does not drive our primary conclusions.

\bigskip\hrule\bigskip

\section*{Reviewer \#2 (Remarks to the Author)}

Review for manuscript COMMSBIO-24-6821-T

\subsection*{Summary of research}

\textit{The manuscript describes attempts to model and forecast microbial abundances and other parameters across terrestrial NEON sites using site specific environmental parameters and current/past microbial community state. The researchers find that bacterial forecasting is more accurate than fungal forecasting and that forecasting is more tractable for certain groupings of taxa (either different taxonomic levels or different functional groups). Despite this, the ability of forecasts to predict variations in observations was still only around 10\% (although this varied depending on how the data was grouped and how the models were run). Beta distributions were used for regression modeling because Dirichlet distributions were computationally intractable.}

\subsection*{Contribution}

\textit{This is a valuable contribution to the literature. Ecological forecasting, while it has gained attention in recent years, is still nascent and difficult to do well. This is particularly true for complex microbial communities that may have very high variability even at local scales. The core finding that bacterial abundances are easier to predict than fungi and the various grouping levels that are most effective are valuable findings for others looking to include microbial taxa in forecast models.}

\subsection*{Overall comments}

\textit{Overall I think this is an important manuscript that makes a novel and valuable contribution to the ecological forecasting and microbial macroecology literature. I think that it would benefit from some revision to make sure the intended meaning is clear and to make it easier for the readership of a broad-scope journal like Communications Biology to understand what was done and why. This includes adding some clarifications throughout in the text, more detailed figure captions, and some cleaning of the code base prior to publication.}

\textit{Red/green plot color issues will make this difficult for color-blind people. There are lots of accessible options out there that would be easy to swap in.}

\textit{The data are all freely available and all code seems to be available as well, but additional organization of the repository and the scripts might be warranted before publication to make it more interpretable/reproducible (script commenting is minimal, contain sections of code are commented out, it's not clear what order things need to be be executed in, setwd() is used in the script, the README doesn't render properly because it has an Rmd extension instead of md, there is a nested microbialForecast directory that looks to be it's own package but that's not included in the README, etc.).}

Thank you for noting these issues! We thank the reviewer for these detailed observations about clarity and reproducibility. Regarding the overall manuscript clarity: we have improved the intended meaning of key results in lines \ref{rev-5}, added rationale for methodological choices in the Methods at lines \ref{rev-6}, and expanded the captions for Figures % TODO: fill in figure numbers
. Regarding the code repository: we have reorganized the GitHub repository so that scripts are clearly commented, the execution order is documented in the README (now rendered as a .md file), setwd() calls have been removed in favor of relative paths, commented-out code sections have been cleaned up, and the nested microbialForecast package directory is now described in the README.

On the issue of accessible figures: we greatly appreciate this note and want our figures to be accessible to the broadest audience possible; in most of our figures, color is present as an ancillary distinction but not essential for reader understanding (e.g. the same variable is splitting by panel in addition to color, and the colors are just used to quicker connection to other figures in the manuscript in supplement).

We tested the main color scheme of our figures (saturated vermillion and bluish-green, as recommended by Okabe \& Ito, \url{https://jfly.uni-koeln.de/color/}) using the Colorblind image tester (\url{https://bioapps.byu.edu/colorblind_image_tester}) and found the resulting colors distinguishable.

Upon re-reviewing the colors for each figure, we have adjusted the color saturation and label visibility for Figure S7, and chosen a different color scheme for Figure S1. If we learn that any specific figures or color palettes remain inaccessible, we will update the material accordingly.

Primary color palette and simulated visibility with 80\% deuteranopia:
% TODO: Include colorblind palette comparison images

\subsection*{Detailed comments}

\textit{Line 47: I'd use another word instead of `domains' here, since fungi are not a domain.}

We changed this word to ``kingdoms'' for consistency with rest of text (line \ref{rev-8}).

\bigskip

\textit{Line 52: This seems like an overly bold statement. In some sites and situations, the change in community composition may be extreme, but in other sites and situations, where seasonal differences are more minor, it may be less so.}

We revised this statement from ``about as dramatic'' to ``potentially as dramatic'' (lines \ref{rev-9}).

\bigskip

\textit{Line 119: Fig 2C refers to fungi, not bacteria, and it looks like from the plot that class is not significantly different from genus, family, or order, so I am not sure this claim is supported.}

We updated this text in the manuscript to clarify that the significant difference in bacterial taxonomic ranks was observed between the class and phylum ranks (lines \ref{rev-10}).

\bigskip

\textit{Line 238: You could test this --- are higher abundance genera (for example) also better predicted?}

We conducted this test and saw non-significant trends in a similar direction for bacteria (R$^2$ = .22, p = .15, n = 11). We now report this on lines \ref{rev-11a}. We did not have enough statistical power to report tests from converged fungal genus-level models (n=4). We report this on lines \ref{rev-11b}.

\bigskip

\textit{Line 249: I believe sodium is considered a micronutrient.}

Removed the phrase ``and macronutrients'' --- instead grouping as a cation, along with potassium and calcium (lines \ref{rev-13}).

\bigskip

\textit{Line 280: How was this cutoff chosen for 236 taxa? What about these 40 functional groups? Even if explained in the cited manuscript, it is very helpful to give a brief rationale.}

The total of 236 taxa was the result of retaining microbial groups with sufficient data for reliable model fitting: specifically, we required a minimum prevalence of 20\% non-zero observations per group. Requiring more observations would have excluded ecologically important but less commonly observed groups; a lower cutoff would have required extreme zero-inflation of models. We have added this justification at lines \ref{rev-14} of the manuscript. Functional groups were chosen for relevance to soil ecological and biogeochemical processes.

\bigskip

\textit{Line 354: I believe it is DADA2 not dada2. That's how it's referred to in the manuscript describing it and on the website.}

We have corrected ``dada2'' to ``DADA2'' throughout the text (see line \ref{rev-dada2}).

\bigskip

\textit{Line 357: Bacterial table size description has number of reads not number of ESVs. I assume this is a typo.}

The description of the bacterial table size has now been corrected to report the number of ESVs rather than the number of reads (line \ref{rev-dada2}).

\bigskip

\textit{Line 368: Are you concerned that since only assigned taxa were put into function groups, that this introduces bias into your analysis? I might add some text here explaining why this is not a problem for your conclusions.}

This is a fair concern. Functional group assignments were limited to classified taxa, meaning unclassified reads are excluded from the functional group relative abundances. This introduces a potential bias if unclassified taxa are systematically enriched in particular functional categories. However, in our dataset, unclassified reads were distributed broadly across the phylogeny rather than concentrated in particular clades, suggesting that their exclusion does not systematically inflate or deflate any particular functional group. We have added a sentence to this effect at line \ref{rev-15}.

\bigskip

\textit{Line 381: I always get a bit nervous when people assign function groups in this way. Ecological function is not inherent to a taxon, but is a product of that taxon and its environmental and host context.}

We agree that this is an important limitation. Guild assignments based on taxonomy assume functional consistency within a taxon regardless of environmental context, which is unlikely to hold universally. As the reviewer notes, the ecological flexibility of microbes --- including context-dependent shifts in function --- means that taxonomically-defined guilds are an imperfect proxy for actual functional roles. We have added text to the Discussion (lines \ref{rev-16}) explicitly acknowledging that forecast skill for functional guilds may be partially constrained by the ecological flexibility of the member taxa, and that future work using trait-based or metagenomics-derived functional assignments would help address this limitation.

\bigskip

\textit{Line 444: How is this handled for evergreen sites? E.g. tropics? How are the effects of winter vs summer precip patterns accounted for (e.g. Mediterranean-type rainfall patterns vs monsoonal)?}

Our seasonal models do not impose a temperate seasonal template on all sites. Rather, seasonality is estimated empirically at each site using trigonometric transformations of sampling month, and peak microbial abundance is linked to plant phenophases derived from the MODIS Land Cover Dynamics product (MCD12Q2), which uses remote-sensing data to estimate site-specific greenup, peak greenness, and dormancy transitions. For evergreen sites, this approach accommodates the phenological cycles that characterize those systems: if no strong seasonal pattern is present in the data, the amplitude of the estimated seasonal term will be near zero, and the model will rely more heavily on environmental predictors. We acknowledge, however, that precipitation seasonality (e.g., winter vs. summer rainfall regimes) is not explicitly represented as a predictor in our models --- though soil moisture is included as a site-monthly mean, and precipitation is included as a site-level mean. We now note this explicitly at lines \ref{rev-17}.

\bigskip

\textit{Line 491: Why MAT and MAP vs min or max annual temp, etc? Presumably many other metrics could be calculated.}

We used variables previously identified as explaining major sources of spatial variation in the NEON microbial community dataset (Averill et al. 2020). At this spatiotemporal scale we explored in this manuscript, the mean, min, and max temperatures are highly correlated (Supplementary Fig/Table % TODO: fill in supplementary number
). While a multitude of additional climate metrics exists, testing them without specific hypotheses about why each metric would be better than MAT and MAP would be both computationally demanding and prone to overfitting.

\bigskip

\textit{Ref 32: Title is all caps.}

\textit{Ref 52: Title is all caps.}

\textit{Ref 77: Looks like it has some remnant LaTeX formatting for acronyms.}

\textit{Ref 90: All caps.}

\textit{Ref 91: All caps.}

\textit{Figure 1: Y-axis should be `Relative Abundance'}

All typographic and formatting issues in the reference list (Refs 32, 52, 77, 90, 91) have been corrected in the revised manuscript. The Figure 1 y-axis label has been updated to ``Relative abundance'' (line \ref{rev-39}).

\bigskip

\textit{Figure 3: Many fungi are both pathogens and endophytes, or endophytes and saprotrophs, depending on context and host. How is that addressed in this classification approach?}

Fungal guild assignments follow FUNGuild conventions, under which taxa can be assigned to multiple guilds (e.g., a taxon may be classified as both a pathogen and an endophyte). Our analysis used primary and secondary assignments, and because assignments are non-mutually exclusive, the summed relative abundances of all guilds can exceed one. We have clarified this in the figure caption and in the Methods at lines \ref{rev-19}.

\bigskip

\textit{Figure 4: Why is the forecast for plot median but the observation it's compared to for the plot mean?}

The figure caption has been updated at lines \ref{rev-20a} to clarify that forecast values represent the posterior median (used because forecast distributions for low-abundance taxa are right-skewed due to zero-inflation), while observed values are plot means. We have also added a note about this distinction in the Supplementary Methods at lines \ref{rev-21}.

\bigskip

\textit{Figure 5B: Is this for fungi and bacteria together? Or one or the other? If it is both, why is it aggregated here but not elsewhere?}

The figure caption has been updated at lines \ref{rev-22a} to clarify that Figure 5B aggregates data across bacteria and fungi, and to explain that this aggregation is valid because the response variable (forecast skill) was normalized by group abundance and predictor variables were normalized within kingdom prior to pooling. We have also noted this explicitly in the Supplementary Methods at lines \ref{rev-22b}.

\bigskip

\textit{Fig S1: Red/green color issues.}

Figure S1 has been updated to use a colorblind-accessible palette, verified using the BYU colorblind image tester.

\bigskip

\textit{Fig S3: X axis label cut off for `Overall predictability'}

The x-axis label for `Overall predictability' in Figure S3 was truncated in the previous version; this has been corrected in the revised figure.

\bigskip

\textit{Fig S3: This is for bacteria? Or all taxa? I assume bacteria because it uses phylogenetic information, but this should be in the caption.}

The caption for Figure S3 has been updated to clarify that this figure shows results for bacteria only, as the phylogenetic information used to construct it is specific to the bacterial dataset.

\bigskip

\textit{Fig S4: X-axis --- I think it's best to choose another word besides `Kingdom' here. Perhaps `Taxonomic Group'? For the bacteria, which type of function group assignment was used here? Weren't there multiple approaches you used? Same question for the taxonomic grouping --- which groups were used? Or were all groups aggregated? It's unclear from the current caption.}

\bigskip

\textit{Fig S6: I am not sure the linear fits from geom\_smooth are appropriate here. The residuals are very definitely not normally distributed.}

\bigskip

\textit{Fig S7: Copiotroph plot has growing season labels cut off.}

The copiotroph plot growing season labels have been corrected in the revised figure (Figure S7).

\bigskip

\textit{Fig S8: Were basic model assumptions met for these linear fits? It seems like subplots A and B for bacteria are particularly off.}

\bigskip

\textit{Fig S9: Are model assumptions met for the linear fits here?}

For Figure S8a--b, where the key finding (flat relationship for fungi, positive relationship for bacteria) is robust to violations of normality, we have retained the fit but added a note in the caption that results should be interpreted qualitatively rather than as formal inference.

\bigskip

\textit{Fig S10: Generally it's helpful for the reader if the take-home is included in the caption. I am assuming here the point is that for these groups, higher seasonal shifts lead to lower error in predictions?}

We have added a take-home statement to each Supplementary Figure caption.

\bigskip

\textit{Fig S15: Have you explored the effects of differing gap fill data sources on prediction capability?}

No, we have not tried to test the effects of differing gap fill data sources on model predictions. This gap-fill analysis was initially developed for model fitting with NEON's provisional sequencing data (2013--2015), which was excluded from the final models (2016--2020). The final dataset includes almost no data gap-filled from SCAN or SMOS, preferring SMAP L4, which was the most closely correlated with the NEON measurements (Ayres et al. 2024).

\bigskip

\textit{Fig S16: If you right-shifted your tick mark labels on the x axis, it would be easier to read.}

The x-axis tick labels in Figure S16 have been right-shifted for legibility in the revised figure.

\bigskip

\textit{Supp refs: Same comments about all caps references needing to be fixed.}

All-caps formatting in the supplementary reference list has been corrected throughout.

\bigskip

Signed: Naupaka Zimmerman

\bigskip\hrule\bigskip

\section*{Reviewer \#3 (Remarks to the Author)}

\textit{This manuscript attempts to forecast the soil microbiome using an extensive spatiotemporal dataset across various ecosystems of the US. The major findings are that the relative abundance of microbial groupings at higher levels of taxonomic classification or functional guilds can be predicted more accurately than lower levels of classification, and that bacteria can be predicted more accurately than fungi, yet fungi can be predicted further in time than bacteria. While I appreciate the ambitious scope of the study and believe the manuscript is well written, I do however have several conceptual issues listed here:}

\textit{The manuscript is written to sound much more far reaching and universal than it actually is. Firstly, Earth and continental scales are often mentioned (including in the title and abstract), which are inappropriate and inaccurate considering the dataset is from the national (US) level. Secondly, it is unclear if it is actually forecasting or just predicting already collected data in a dataset. Furthermore, abundance is often used instead of relative abundance, and the weaknesses of using relative abundances and guild assignments based on metabarcoding data is not discussed. For example, why would we want to forecast the relative abundance of microbial guilds and taxa instead of trying to forecast activity or actual processes? What would be the importance of forecasting the relative abundance of microbial guilds for any given application? I particularly take issue with using the relative abundances at the kingdom level, this should be accompanied by some sort of proxy measure of abundance such as qPCR or PLFAs, which would be more informative.}

We thank the reviewer for these critiques and for the opportunity to clarify the scope and motivation of our work.

\begin{itemize}
\item \textbf{Geographic scope:} ``Continental'' and ``Earth-scale'' language has been replaced throughout with ``national-scale'' or ``across the United States'' (title, abstract lines \ref{rev-24}, main text lines \ref{rev-25}). We have retained references to the breadth of environmental gradients captured (spanning % TODO: fill in numbers
 ecoclimatic domains and % TODO: fill in numbers
 biomes) to convey the ecological scope while being accurate about geographic boundaries.

\item \textbf{Relative abundance:} ``Relative abundance'' is now used consistently throughout the manuscript (e.g., lines \ref{rev-27a}, \ref{rev-27b}, \ref{rev-27c}). Regarding the value of forecasting relative abundance specifically: microbial community composition is directly linked to key ecosystem processes including decomposition, nitrogen cycling, and carbon storage (e.g., % TODO: add citation
), and relative abundance is the most widely available community-level metric from large-scale sequencing surveys. Beyond applied motivations, understanding whether and at what scales microbial community composition is predictable is a fundamental question in community ecology. We have added text to the manuscript at lines \ref{rev-29} making these motivations explicit.

\item \textbf{qPCR/PLFA:} We agree that absolute abundance measures such as qPCR or PLFA would complement our relative abundance data, particularly for kingdom-level analyses. We now discuss this explicitly at lines \ref{rev-30}, acknowledging that relative abundance forecasts cannot capture absolute biomass dynamics and that future work incorporating qPCR or PLFA data would strengthen conclusions about total microbial activity.
\end{itemize}

\bigskip

\textit{It is unclear what type of ecosystems were surveyed and why particular explanatory variables were chosen. For example, what kinds of vegetation were present, why wasn't vegetation composition or some sort of plant phylogenetic variable used instead of ectomycorrhizal tree basal area? Were there grasslands, if so why EM basal area? Were there coniferous forests? If so this would be important to include in the models. Why percent carbon, and not C:N, or \% N, I imagine these would improve the models? Also why not a spatial variable like the main PCNM axis or even coordinates to better explain microbial relative abundances, due to the importance of distance decay in similarity? I understand that site is included as a random factor, but I still think the spatial effect is understated or under-examined and would improve the models.}

Our predictor variables were selected in a prior analysis (Averill et al. 2022) that systematically evaluated soil and vegetation properties for their ability to explain spatial variation in microbial community composition across NEON sites while minimizing collinearity and overfitting. Regarding specific variables: (1) EM tree basal area was included because it is one of the strongest predictors of fungal community composition at continental scales and is present at all NEON terrestrial sites, including grasslands (where its value is near zero, providing meaningful contrast); (2) soil carbon percentage was selected over C:N because it explained more variance in the prior analysis and C:N was collinear with other predictors; (3) we did not include PCNM axes or raw spatial coordinates as fixed effects because we found no residual spatial autocorrelation in model residuals after fitting the site random effects to environmental predictors (Moran's I = % TODO: fill in statistical values
, p = % TODO: fill in statistical values
), indicating that the spatial signal was already captured. We have added justification for these choices at lines \ref{rev-33} of the revised manuscript.

\subsection*{Line specific comments}

\textit{Title: replace continental to national, the US is not a continent}

\textit{L3--4 Abstract: It is beyond the scope of this study the Earth's microbiome, and it comes across as a bit grandiose to mention it. Why not just mention the soil microbiome across different regions. Also this study is not at the continental-scale, please change to reflect the national scale instead.}

\bigskip

\textit{L5--6 Abstract: From my understanding the data is relative abundances, not abundances. I have comments about the idea of fungal and bacterial relative abundances further below.}

We have systematically revised the manuscript to use ``relative abundance'' consistently throughout (e.g., lines \ref{rev-27a}, \ref{rev-27b}, \ref{rev-27c}, \ref{rev-44}). We acknowledge this precision is critical for interpreting results from amplicon sequencing data and have ensured consistency across all figures, tables, and text.

\bigskip

\textit{L11 Abstract: Once again, overkill to scale this study to be relevant for Earth.}

\bigskip

\textit{L65: I would change powered to a more appropriate word like facilitated or enabled}

Revised accordingly: ``powered'' replaced with ``facilitated'' in line \ref{rev-35}.

\bigskip

\textit{L66: It would be useful to mention here briefly what sort of ecosystems are included in this network, is it just forests?}

We have revised line \ref{rev-36} to note that the NEON network spans % TODO: fill in numbers
 ecoclimatic domains, encompassing the majority of terrestrial biome types present in the United States, including temperate and boreal forests, grasslands, shrublands, wetlands, and tundra.

\bigskip

\textit{L85--87: 10\% of variation explained by forecasts seems like a very low cutoff for considering them reasonable, it seems that this would be much better than a semi informed guess, or even a random guess.}

This value of 10\% as a threshold for ``reasonable'' model skill was chosen pragmatically to distinguish models that capture meaningful temporal and spatial signal from those that are essentially uninformative relative to a null mean model. In a system as variable as soil microbial communities, an R$^2$ of 10\% represents a non-trivial departure from chance. We now report the full distribution of R$^2$ values across groups in Supplementary Table % TODO: fill in supplementary number
 so readers can apply their own criteria.

\bigskip

\textit{Figure 1: The y-axis should be relative abundance. The caption}

The y-axis label for Figure 1 has been updated to ``Relative abundance'' (line \ref{rev-39} of the figure caption).

\bigskip

\textit{Figure 2: A description is lacking for subfigure f (which should be e), it is unclear what is being predicted within the categories bacteria and fungi in this case.}

Thank you for identifying this --- the Figure 2 caption has been corrected (line \ref{rev-fig2fix}).

\bigskip

\textit{L116--120: It is obvious that larger groupings are easier to predict than more precise groupings, this is well established.}

While it may seem intuitive that broader taxonomic groupings should be easier to predict, the ecological literature has in fact suggested the opposite: that taxonomic aggregation obscures ecologically meaningful variation, dilutes the signal of biotic interactions, and may reduce the explanatory power of environmental predictors (e.g., % TODO: add citation
). This hypothesis had not previously been tested using true forward forecasts --- i.e., predictions made at new time points and new locations using models fitted on held-out data. To our knowledge, this was the first study to conduct such an unbiased test at national scale for soil microbial communities. The empirical finding that aggregation improves forecast skill is therefore not obvious but rather a novel and testable result with practical implications for future forecasting efforts.

\bigskip

\textit{L141--142: Why didn't you use partially-explicit modeling approaches? This would also greatly improve your fungal predictions.}

Our models are fully data-driven state-space dynamic linear models, in which temporal dynamics are captured by the lag-1 autoregressive term and environmental predictors, with all parameters estimated from the data. We interpret ``partially-explicit'' modeling as approaches that fix some parameters based on prior mechanistic knowledge --- for example, embedding known microbial growth or stoichiometric relationships within an otherwise statistical framework. We chose not to pursue this approach for multiple reasons: first, the mechanistic parameters necessary to inform such constraints (e.g., taxon-specific growth rates, carbon use efficiencies, interaction coefficients) are not known at the taxonomic and geographic breadth of this dataset. Given this, imposing partially-informed structure with varying levels of certainty across groups would introduce inconsistency in comparisons across groups.

\bigskip

\textit{L145--146: Yes for well defined and clear ecological functions, but such categories do not capture the ecological flexibility of microbes.}

We agree with the reviewer's point. The functional guild assignments used here treat each taxon as belonging to a fixed functional category, but in reality most microbes perform a shifting repertoire of functions depending on resource availability, redox conditions, and community context. A nitrogen-cycling bacterium, for instance, may switch between nitrification and denitrification pathways depending on oxygen availability, and many fungi shift between saprotrophic and symbiotic strategies depending on host and substrate. Our guild assignments therefore capture a useful average tendency rather than a fixed functional identity. We have added text to the Discussion at lines \ref{rev-41} acknowledging that the forecast skill of functional guilds is likely partially constrained by this context-dependence, and that future approaches incorporating dynamic functional measurements --- such as metatranscriptomics or enzyme assays --- would better capture the realized functional roles of microbial communities at any given point in time.

\bigskip

\textit{Figure 4: It seems that these results might be an artifact of amplicon sequencing, how have you determined relative abundances of bacteria and fungi? Is it the proportion of total reads assigned to bacteria and fungi from 16S and ITS sequencing? Does it include unassigned reads? If so this seems like a very flimsy measure of abundance that is very sensitive to choice of primers and PCR conditions. I imagine more reads of the 16S dataset were assigned to bacteria, than ITS reads were assigned to fungi, which could explain such differences. Or have I completely misunderstood what these graphs represent?}

Bacterial and fungal relative abundances were determined from separate amplicon sequencing datasets, processed independently. Relative abundances reflect the proportion of reads assigned to each taxon within its respective dataset; unassigned reads were excluded following standard practice in the amplicon sequencing literature (e.g., % TODO: add citation
), as their taxonomic identity is uncertain. Crucially, because bacterial and fungal datasets were sequenced and analyzed separately, differences in the proportion of assigned reads between the two kingdoms cannot explain differences in forecast skill between them.

We agree that the choice of primers will inevitably determine the taxonomic composition of the data produced. NEON updated their primer sets in 2015, which could introduce a systematic shift in detected community composition. We addressed this directly by including a binary covariate (``legacy'') in every model to capture community-level shifts associated with this primer change, isolating it from the ecological dynamics we are forecasting. This covariate is defined in the model notation (Appendix S3, Notation section) as ``a binary indicator for early sequencing timepoints that used different amplification primers, creating a systematic bias in relative abundances across the time series,'' and its prior is specified in the Summary of Prior Distributions (Appendix S3).

\bigskip

\textit{L259: I would change fundamentally to theoretically, or just remove all together and leave just possible.}

Revised accordingly --- ``fundamentally'' has been removed, now reading ``forecasting the soil microbiome in time is now possible'' (line \ref{rev-29}).

\bigskip

\textit{L276--293: Okay so I am yet to come across if it is explicitly stated that you did the analysis of the calibration period before the `forecast' data was already collected, were these predictions made after the fact, in that case they are not forecasts and that term should be removed from the manuscript unless I am mistaken? Wouldn't it only be model predictions in this case rather than forecasts?}

We appreciate the opportunity to clarify our use of the term ``forecast.'' The use of ``forecasts'' is from the model's point of view (typical in climate science, e.g. \url{https://applets.kcvs.ca/ClimateModelHindcasting/}), and these forecasting models (which were evaluated using a time series not included in the calibration period of 2016--2018) were used for the majority of results described in this paper.

Evaluations of the models with all data assimilated (2016--2020) could be considered hindcasts, not true forecasts, but this is primarily a data latency problem: if more recent sequencing data had been released at the time of analysis, the forecasts would have extended into the future using weather forecasts, but the COVID-19 pandemic created a gap in data availability during this research period. ``Hindcast'' is a niche term typically used within climate science communities, and we intend for this paper to be useful to the broader ecological forecasting community, hence the use of ``forecasts'' instead of ``hindcasts'' in the title.

We use the term forecast rather than prediction due to its strong connotations with probabilistic time-series analysis based on prior trends, whereas ``prediction'' often refers to in-sample predicted values. This context has been added in the newest revision.

\bigskip

\textit{L280: This manuscript often uses abundances and relative abundances interchangeably, which is misleading.}

As described in our response to the general comments above, ``relative abundance'' is now used consistently throughout the manuscript, including in the Methods section at lines \ref{rev-44}. We have also added explicit clarification that no direct measure of absolute microbial abundance was obtained, and that all modeling was performed on amplicon-derived relative abundances.

\bigskip

\textit{L302--303: Why these predictor variables? I see no justification? For example why soil C \% and not C:N or N \%? And also why EM trees, and not other vegetation properties?}

Soil and vegetation properties were chosen in a previous study that attempted to predict relative abundances of microbial groups only in new locations (Averill et al. 2022). This analysis attempted to select predictors with explanatory power across many microbial groups while minimizing correlation between predictors and avoiding overfitting / data dredging. We now explain this in the manuscript (lines \ref{rev-45}).

\bigskip

\textit{L314: why not use actual spatial measures such as coordinates and distance from one another?}

We feel confident that the answers we arrive at in this analysis are not dependent on unexplained spatial patterns, because there was no residual spatial autocorrelation detected between NEON sites after fitting the random effects with spatial predictors. We chose not to create a fully spatially-explicit model for a number of reasons, which we have attempted to clarify in the newest revision. One reason was computational tractability: our attempts to create the first continuous time-series for hundreds for microbial groups required substantial resources to evaluate choices in model structure and parameterization for an understudied ecological system. We consider a fully spatial model to be a future goal, after resolving initial broad questions about data patterns.

\bigskip

\textit{L326--330: Please give a brief overview of the types of vegetation types sampled, was it grasslands, forests? Coniferous forests, broadleaf forests? This will help reviewers assess the appropriateness of using EM tree \% as an explanatory variable.}

A complete list of NEON sites included in our analysis, including vegetation type, is provided in Supplementary Table % TODO: fill in supplementary number
. The sites span a broad range of vegetation types, including temperate deciduous and coniferous forests, grasslands, shrublands, wetlands, and tundra. The full scope of the NEON observatory and its site selection criteria are described in NEON design publications (% TODO: add citation
). The use of EM tree basal area as a predictor is appropriate across this range because its value is zero at non-forested sites, providing meaningful ecological contrast without requiring site-type-specific models.

\bigskip

\textit{L338: Until further?}

Revised accordingly --- corrected to ``until further microbial and chemical analyses'' (line \ref{rev-48}).

\bigskip

\textit{L339: Did you account for sampled horizon type in the analyses? This is very important.}

We decided to only model the dominant horizon type per site, which was generally the organic horizon, with the mineral horizon only used for sites entirely lacking a labeled organic horizon. Variation between horizons was mostly explained by soil physicochemical properties during exploratory data analysis, and residual variation would be captured by the site effects. This decision was additionally informed by known data quality issues with horizon classification in early NEON releases, where seasonal technician assignments occasionally required correction. This explanation is now included in the Methods (lines \ref{rev-48}).

\bigskip

\textit{L340--350: Nowhere here have I come across a method for measuring abundance, which is often mentioned in the manuscript.}

``Relative abundance'' is now used consistently throughout the manuscript, including in the Methods section at lines \ref{rev-44}. We have also added explicit clarification that no direct measure of absolute microbial abundance was obtained, and that all modeling was performed on amplicon-derived relative abundances.

\bigskip

\textit{L344: Which subregions of 16S?}

NEON uses 16S primers that amplify the V3--V4 region of DNA coding for ribosomal RNA. We have specified in newest revision on lines \ref{rev-50}.

\subsection*{Works cited}

E. Ayres, R. H. Reichle, A. Colliander, M. H. Cosh and L. Smith, ``Validation of Remotely Sensed and Modeled Soil Moisture at Forested and Unforested NEON Sites,'' in \textit{IEEE Journal of Selected Topics in Applied Earth Observations and Remote Sensing}, vol. 17, pp. 14248--14264, 2024, doi: 10.1109/JSTARS.2024.3430928.

\end{document}
